% =============================================================================
%  Notas de clase — Sesión 13: Torres de Hanoi — Solución Recursiva
%  Programación Avanzada — Universidad Panamericana
%  Compilar: pdflatex sesion13_hanoi_recursivo.tex
% =============================================================================
\documentclass[12pt, oneside]{article}
\usepackage{progavanzada}

\begin{document}

\portada{Sesión 13}{Torres de Hanoi: Solución Recursiva}{2026}

\tableofcontents
\newpage

% =============================================================================
\section{El concurso de diseño}
% =============================================================================

Antes de ver el algoritmo de hoy, el grupo elige el mejor diseño de la
sesión 12.  Las reglas del concurso son simples:

\begin{itemize}
  \item Cada persona vota por \textbf{exactamente tres} equipos.
  \item Los criterios son: diseño visual atractivo \textbf{y} código legible.
        Los dos cuentan.  Un diseño bonito con código ilegible no gana.
  \item El equipo con más votos comparte su código en pantalla.
  \item Ese código se usará como punto de partida para la actividad de hoy.
\end{itemize}

\begin{consejo}
  Cuando estés votando, pregúntate: ¿entiendo lo que hace cada función con
  solo leer el nombre? ¿Las torres se ven bien en todos los estados?
  Un voto bien razonado pesa más que uno al azar.
\end{consejo}

% =============================================================================
\section{¿Se puede resolver Hanoi en seis líneas?}
% =============================================================================

La sesión 12 te tomó 50 minutos de trabajo en equipo para mover los discos
a mano.  Hoy verás que \textbf{la solución recursiva tiene seis líneas de código}.

¿Cómo es posible?  Porque la descripción recursiva del problema \textit{ya es}
el algoritmo.  Recuerda la lógica de la sesión 12:

\begin{quote}
  Para mover $n$ discos de A a C usando B como auxiliar:\\
  1. Mueve los $n-1$ discos superiores de A a B (usando C como auxiliar).\\
  2. Mueve el disco $n$ (el más grande) de A a C.\\
  3. Mueve los $n-1$ discos de B a C (usando A como auxiliar).
\end{quote}

Eso en C++ se escribe exactamente así:

\begin{codigo}
void hanoi(int n, char origen, char destino, char aux,
           vector<int>& torA, vector<int>& torB, vector<int>& torC) {
    if (n == 0) return;
    hanoi(n-1, origen, aux, destino, torA, torB, torC);
    moverDisco(origen, destino, torA, torB, torC);
    hanoi(n-1, aux, destino, origen, torA, torB, torC);
}
\end{codigo}

El caso base es $n = 0$: no hay nada que mover, se regresa de inmediato.
Todo lo demás es consecuencia de los tres pasos descritos arriba.

\begin{recuerda}
  La recursión no es magia: es la misma lógica que aplicaste a mano en la
  sesión 12, traducida directamente a código.  La computadora simplemente
  la ejecuta $2^n - 1$ veces sin cansarse.
\end{recuerda}

% =============================================================================
\section{El árbol de llamadas}
% =============================================================================

Cada llamada a \texttt{hanoi(n, ...)} genera \textbf{dos llamadas}
a \texttt{hanoi(n-1, ...)}.  El árbol de recursión para $n = 3$ tiene
esta forma:

\begin{verbatim}
                     hanoi(3, A, C, B)
                    /                 \
         hanoi(2, A, B, C)        hanoi(2, B, C, A)
          /          \              /          \
  hanoi(1,A,C,B) hanoi(1,C,B,A)  hanoi(1,B,A,C) hanoi(1,A,C,B)
    /    \          /    \          /    \          /    \
  h(0)  h(0)    h(0)  h(0)      h(0)  h(0)      h(0)  h(0)
\end{verbatim}

Las hojas (\texttt{h(0)}) son los casos base.  Los nodos internos son las
llamadas que hacen trabajo real (\texttt{moverDisco}).

Contando los nodos internos del árbol:

\begin{center}
\begin{tabular}{cccc}
  \toprule
  \textbf{$n$} & \textbf{Hojas} & \textbf{Nodos internos} & \textbf{Movimientos} \\
  \midrule
  1 & 2  & 1   & 1 \\
  2 & 4  & 3   & 3 \\
  3 & 8  & 7   & 7 \\
  4 & 16 & 15  & 15 \\
  $n$ & $2^n$ & $2^n - 1$ & $2^n - 1$ \\
  \bottomrule
\end{tabular}
\end{center}

Cada nodo interno corresponde exactamente a un movimiento.
Por eso $M(n) = 2^n - 1$: es el número de nodos internos de un árbol
binario completo de profundidad $n$.

% =============================================================================
\section{El call stack durante la ejecución}
% =============================================================================

La profundidad del call stack en cualquier momento es \textbf{como máximo $n$},
porque cada llamada espera a que su primera sub-llamada ($n-1$) termine antes
de continuar.  Nunca hay más de $n$ frames activos simultáneamente:

\begin{verbatim}
Momento en que se mueve el disco 3 (paso 4 de 7):

  hanoi(3, A→C, B)   ← frame más antiguo, esperando
    hanoi(2, A→B, C) ← ya terminó
    [moverDisco A→C] ← ← línea activa aquí
    hanoi(2, B→C, A) ← aún no llamada
\end{verbatim}

La profundidad máxima es $n$, no $2^n$.
Para $n = 64$ el stack tendría solo 64 frames; lo que explota es el
\textit{tiempo}, no la memoria.

\begin{tecnico}[Complejidad temporal y espacial]
  La función \texttt{hanoi} tiene:
  \begin{itemize}
    \item \textbf{Complejidad temporal:} $O(2^n)$ — el número de llamadas totales
          crece exponencialmente con $n$.
    \item \textbf{Complejidad espacial:} $O(n)$ — la profundidad máxima del
          call stack en cualquier momento es $n$.
  \end{itemize}
  Muchos algoritmos recursivos tienen complejidad espacial proporcional a la
  \textit{profundidad del árbol}, no al número total de llamadas.
\end{tecnico}

% =============================================================================
\section{Verificación: contando movimientos}
% =============================================================================

Una forma de verificar que el algoritmo es correcto es contar cuántas veces
se llama \texttt{moverDisco}.  Si el programa mueve $n$ discos, el contador
debe ser exactamente $2^n - 1$.

\begin{codigo}
int movimientos = 0;

void hanoi(int n, char origen, char destino, char aux,
           vector<int>& torA, vector<int>& torB, vector<int>& torC) {
    if (n == 0) return;
    hanoi(n-1, origen, aux, destino, torA, torB, torC);
    moverDisco(origen, destino, torA, torB, torC);
    movimientos++;
    hanoi(n-1, aux, destino, origen, torA, torB, torC);
}
\end{codigo}

Al terminar, \texttt{movimientos == pow(2, n) - 1}.

% =============================================================================
\section{La actividad de hoy}
% =============================================================================

Trabajas con el código ganador del concurso como base.

\subsection{Lo que debes implementar}

\begin{enumerate}
  \item \textbf{La función \texttt{hanoi}} que resuelve el problema automáticamente,
        usando \texttt{moverDisco} del código base para actualizar el estado
        de las torres.

  \item \textbf{Imprimir las torres después de cada movimiento} para que se
        pueda seguir la solución visualmente.

  \item \textbf{Contar y mostrar} el número de movimientos al final y
        verificar que sea $2^n - 1$.
\end{enumerate}

\subsection{Criterios de evaluación}

\begin{center}
\begin{tabular}{lp{7cm}c}
  \toprule
  \textbf{Criterio} & \textbf{Descripción} & \textbf{Peso} \\
  \midrule
  Corrección      & El programa mueve todos los discos de A a C
                    respetando las reglas en cada paso & 40\% \\
  Conteo          & Al terminar, imprime el número de movimientos
                    y verifica que sea $2^n - 1$ & 20\% \\
  Visualización   & Las torres se muestran (o pueden mostrarse)
                    después de cada movimiento & 20\% \\
  Integración     & La función recursiva usa \texttt{moverDisco}
                    del código base sin reescribirla & 20\% \\
  \bottomrule
  & \textbf{Total} & 100\% \\
  \bottomrule
\end{tabular}
\end{center}

\begin{advertencia}
  Muchas soluciones de Hanoi en internet imprimen directamente los movimientos
  (\texttt{cout << "Mueve de " << A << " a " << C}).
  Eso no es lo que se pide: el objetivo es que la función use
  \texttt{moverDisco} para actualizar el \texttt{vector}, de modo que las
  torres se puedan dibujar con el diseño elegante de la sesión 12.
\end{advertencia}

% =============================================================================
\section{Ejercicios}
% =============================================================================

\begin{enumerate}
  \item Traza a mano el árbol de recursión de \texttt{hanoi(2, A, C, B)}.
        Indica en qué orden se hacen los movimientos y qué disco se mueve
        en cada uno.

  \item ¿Cuántas veces se llama \texttt{hanoi} en total (incluyendo las
        llamadas con $n=0$) cuando se invoca \texttt{hanoi(4, ...)}?
        \textit{(No olvides las hojas del árbol.)}

  \item Modifica \texttt{hanoi} para que reciba un parámetro
        \texttt{int nivel} (inicialmente 0) y muestre una indentación
        proporcional al nivel antes de cada movimiento, igual que
        hiciste con \texttt{sumaTraza} en la sesión 11.

  \item ¿Es posible implementar Hanoi de forma iterativa (sin recursión)?
        Si crees que sí, describe en palabras cómo sería la lógica.
        Si crees que no, justifica por qué.

  \item \textbf{Desafío:}
        Modifica el programa para que el usuario pueda elegir
        el número de discos $n$ al inicio.
        Añade un límite: si $n > 20$, muestra solo el conteo de movimientos
        sin dibujar las torres (para no inundar la terminal).
\end{enumerate}

% =============================================================================
\section{Resumen}
% =============================================================================

\begin{recuerda}
\begin{itemize}
  \item La solución recursiva de Hanoi sigue exactamente los mismos tres pasos
        que aplicaste a mano en la sesión 12.
  \item El código tiene \textbf{una} condición base ($n = 0$) y
        \textbf{tres} líneas de trabajo: sub-llamada, mover, sub-llamada.
  \item El árbol de recursión tiene $2^n - 1$ nodos internos, uno por movimiento.
  \item La profundidad del call stack es $O(n)$; el tiempo es $O(2^n)$.
  \item El número de movimientos se puede verificar: siempre debe ser
        exactamente $2^n - 1$.
\end{itemize}
\end{recuerda}

\end{document}
